% !TEX root = ../main.tex
\section{Institutional Warrant and Tiered Reliance}
\label{sec:tiers}

The architecture of Section~\ref{sec:governed-rag} serves a condition I have so far
used without defining. This section defines it and makes the controls proportional
to consequence.

\subsection{Warrant}

Grounding is necessary for reliance and not sufficient. An assertion can rest on
authentic, applicable, authoritative evidence, be traceable to it, and still be
something the institution is not entitled to act on---because nobody checked it, or
because nobody with the standing to authorise reliance did so.

\begin{definition}[Institutional warrant]
\label{def:warrant}
\begin{equation}
  \begin{aligned}
  \Warranted_I(p \mid c) \;\leftrightarrow\;
    &\Resolved(c) \wedge \Grd_S(p \mid c)\\
    &\wedge\ \Verified(V, p, c) \wedge \Authorized(I, p, c).
  \end{aligned}
  \label{eq:warrant}
\end{equation}
\end{definition}

where

\begin{itemize}
  \item $\Resolved(c)$: the material context variables of \eqref{eq:context} are
        explicit rather than assumed or inferred without validation;
  \item $\Grd_S(p \mid c)$: Definition~\ref{def:grounding}---which unpacks to the
        six evidence conditions of \eqref{eq:grounded} plus $\Sensitive$, so
        authenticity, versioning, authority, and applicability are preconditions of
        valid support rather than separate requirements;
  \item $\Verified(V, p, c)$: an appropriate procedure has checked the assertion,
        its evidence relation, and the controls the decision class requires;
  \item $\Authorized(I, p, c)$: the institution permits reliance at the relevant
        risk tier, possibly subject to human review or approval.
\end{itemize}

Three properties of \eqref{eq:warrant} matter for how it is used.

It is a conjunction, not a score. There is no aggregation of the four conjuncts into
a confidence figure, and an assertion satisfying three of them is not
three-quarters warranted. This is the same point as the flow in
Section~\ref{sec:governed-rag}: these are distinct conditions with distinct owners
and distinct failure records.

It is a property of a process, not of a model. Nothing in \eqref{eq:warrant} is a
property of model weights, a benchmark score, or a generated response.
$\Authorized$ in particular is a fact about an institution's governing rules and
delegations, and no improvement to $M$ can supply it. This is why warrant cannot be
procured: a vendor can sell you a component that makes grounding measurable, and
cannot sell you an authorization.

It is revocable and context-relative. \eqref{eq:warrant} is not a claim that $p$ is
certain, or even that $p$ is true. It is the condition under which an institution
may responsibly act on $p$ for a defined purpose, knowing that it may be wrong and
that the record will show what was relied upon and why. That is the standard
ordinary institutional decision-making already meets, and the standard machine
assertion currently does not.

\subsection{Reliance tiers}

Controls should be proportional to the consequence of reliance. The table below is
deliberately compact: its purpose is to show that permissible reliance is set by
consequence rather than by model quality, and that argument does not require a
fully specified control catalogue. Any real programme will need its own tiering,
calibrated to its regulator and its risk appetite.

\begin{table}[ht]
  \centering
  \small
  \begin{tabular}{@{}>{\raggedright\arraybackslash}p{1.1cm}>{\raggedright\arraybackslash}p{3.6cm}>{\raggedright\arraybackslash}p{4.4cm}>{\raggedright\arraybackslash}p{4.4cm}@{}}
    \toprule
    Tier & Use & Required & Invariants engaged \\
    \midrule
    0 & Unwarranted generative assistance: brainstorming, drafting
      & Clear framing as unverified; no reliance claim
      & none \\
    \addlinespace
    1 & Informational assistance
      & Visible uncertainty; no reliance claim
      & I17 \\
    \addlinespace
    2 & Evidence-linked assistance
      & Sources visible; user retains verification duty
      & I8, I9, I12 \\
    \addlinespace
    3 & Governed decision support
      & Context-bound evidence, verification controls, audit trail, authorized
        human reliance
      & I8--I19 \\
    \addlinespace
    4 & Constrained execution
      & Deterministic guardrails, approvals, reversibility, monitoring, full
        records
      & I8--I19 plus action controls \\
    \addlinespace
    5 & Autonomous consequential action
      & Narrowly specified settings, stringent formal controls, continuing
        assurance, explicit authority to act
      & all, plus independent assurance \\
    \bottomrule
  \end{tabular}
  \caption{Reliance tiers. The tier is set by the consequence of reliance, not by
    model quality. Tier~5 is included for completeness and should be approached
    with considerable caution in regulated settings; nothing in this paper argues
    that current systems are ready for it.}
  \label{tab:tiers}
\end{table}

Assignment to a tier needs criteria, and while a full consequence taxonomy is out of
scope, the dimensions are nameable: \emph{reversibility} of the action the assertion
supports; \emph{magnitude} of the exposure if it is wrong; \emph{breadth}, meaning how
many parties are affected and whether they consented; \emph{detectability}, whether an
error would surface before harm; and \emph{availability of a competent human check} in
the time the decision allows. A use case scoring badly on reversibility and
detectability belongs at a higher tier than its nominal risk rating suggests, because
those two dimensions determine whether anything can be done after the fact.

The principle the table exists to enforce is a negative one. Permissible reliance
does not rise because the model got better. It rises when evidence, verification,
constraints, authorization, and accountability rise together and in proportion to
what is at stake. A system that improves its benchmark score has not thereby earned
a higher tier, and a system that has implemented I8--I19 has not thereby earned
Tier~5.

Two boundaries are worth marking explicitly, because they are where programmes
drift. The step from Tier~2 to Tier~3 is the step at which the institution rather
than the user assumes the verification duty, and it requires everything in
Table~\ref{tab:invariants-system}; systems are frequently described as Tier~3 while
implementing Tier~2. And the step from Tier~3 to Tier~4 is the step from assertion
to action, at which reversibility and blast radius become first-order design
questions rather than operational details.
